\SourceLeaf{161}$m, m', m'', m'''$ etc. themselves, then of their ratio (i. e. the ratio of the numbers by which they are expressed, with arbitrary units chosen for each in the heterogeneous case). With ordinary ingenuity one can almost always find at least some means of estimation. Where all knowledge is entirely lacking, it is surely no fault of the method itself that it cannot be of service. One can never dispense with \textls{knowledge of the subject} when applying the method of least squares. It will therefore always depend on \textls{knowledge of the subject} to what extent, in \textls{linear measurements}, the mean error to be expected is to be taken as equal or unequal for shorter and longer distances, and, in the latter case, how it is to be taken to depend on the magnitude. Consider, for example, short distances. If, say, they are measured with a measuring rod 3 metres long, then for a distance of \textls{very few} metres the error will depend entirely or almost entirely on the graduations. But if the rod must be turned end for end many hundreds of times, a quite different source of error comes into play.

\begin{center}
. . . . . . . . . . . .
\end{center}
